\section{Introduction}
\label{sec:intro}

Retrieval-augmented generation (RAG) gives a language model access to text that its parameters do not contain, by retrieving passages and placing them in the prompt \citep{lewis2020rag}. For questions that require combining several facts, as in multi-hop benchmarks such as MuSiQue, 2WikiMultiHopQA and HotpotQA \citep{trivedi2022musique,ho2020twowiki,yang2018hotpotqa}, the passages are typically handed to the model as an unordered list and left for the model to link.

A growing family of methods therefore transforms retrieved evidence before generation. Passages are decomposed into propositions or atomic claims \citep{chen2024densex}, organised into trees of summaries \citep{sarthi2024raptor}, or linked into entity and fact graphs whose nodes, relations and paths are serialised into the prompt \citep{edge2024graphrag,gutierrez2024hipporag,gutierrez2025hipporag2,zhu2025kg2rag,guo2024lightrag,chen2026pathrag}.

Every such transformation can also go wrong. An extractor may omit the sentence that names the answer, and a token budget may cut it because the structured format spends tokens on identifiers and relation lines. Splitting a passage into sentences turns ``It stars Anne Hathaway \dots'' into an item whose subject must be recovered from elsewhere. A relation predictor may add a statement that the source never made: in our experiments, it marked ``The teaser of the film was released in February 2016'' and ``The trailer of the film was released in October 2016'' as contradicting each other, and once shown that line, the model declared the evidence insufficient for a question it otherwise answered correctly.

Benefits and costs therefore arrive together, and the system-level comparisons in the literature cannot separate them. When a graph-based system outperforms or underperforms passage RAG \citep{xiang2026whentographs,han2025ragvsgraphrag,zhou2025indepth}, its retrieval, the content that reaches the model and the way that content is formatted all differ at once.

We argue that the effect of an evidence representation should be measured with the content held fixed, and that the cost of a transformation, measured as the evidence it loses or the information it adds, should be reported separately from its benefit. Therefore, we investigate how the representation of the \emph{same} supporting evidence affects multi-hop answer accuracy, and when transformation errors outweigh the benefits of explicit structure.

To do so, we present the gold supporting evidence of each question in seven representations: passages, sentence items, shuffled items, reasoning steps, graph-ordered items, graph-ordered items with explicit relation lines, and the same with predicted relations. An \emph{information-controlled} protocol verifies, per question, that every representation shows exactly the same facts; a \emph{budget-controlled} protocol gives every representation the same token budget and the same representation-independent ranking of facts, and records which facts each representation can afford. We complement these with retrieved-evidence baselines and controlled perturbations. All design decisions, primary contrasts and statistics were fixed before generation. After an exploratory pilot, we tested the six primary contrasts on a disjoint confirmatory sample of 4{,}976 questions, with sample sizes set from the pilot variance.

\paragraph{Contributions.} (i)~A controlled framework that separates the information a representation carries from its formatting and from the information its structure adds, with per-question verification of fact identity and accounting of evidence loss under a budget. (ii)~A pre-registered test with mostly null results: for Qwen2.5-7B-Instruct, neither graph-derived ordering nor explicit derived relations changed exact match detectably on MuSiQue or 2WikiMultiHopQA, while splitting passages into sentence items lowered 2WikiMultiHopQA exact match by 1.7 points, almost entirely through increased abstention. (iii)~Evidence that noisy structure hurts: predicted relations lowered exact match by four to five points, a loss traced to spurious contradiction relations, and under a token budget structured formats kept 15--20 percentage points fewer facts than passages. (iv)~A release of code, frozen question sets, run manifests and analysis outputs, so that every number can be traced to a run.
